\clearpage
\phantomsection
\addcontentsline{toc}{chapter}{TÓM TẮT}

\begin{center}
% 1. TÊN ĐỀ TÀI Ở TRÊN (Size 14, in hoa, in đậm) theo Mẫu 5
{\fontsize{14pt}{18pt}\selectfont \bfseries \MakeUppercase{\tendetai}\par}
\vspace{0.6cm}

% 2. TÓM TẮT Ở DƯỚI (Size 14, in hoa, in đậm) theo Mẫu 5
{\fontsize{14pt}{18pt}\selectfont \bfseries TÓM TẮT\par}
\vspace{0.8cm}
\end{center}

Trong kỷ nguyên chuyển đổi số và phát triển du lịch thông minh, việc ứng dụng các Mô hình Ngôn ngữ Lớn (Large Language Models -- LLM) nhằm hỗ trợ du khách tra cứu thông tin và tự động thiết kế lịch trình tham quan ngày càng trở nên phổ biến. Mặc dù sở hữu năng lực tạo sinh ngôn ngữ tự nhiên vượt trội, các hệ thống LLM thương mại độc lập cùng các kiến trúc Tăng cường Tạo sinh Truy xuất truyền thống (Naive RAG) thường xuyên bộc lộ hai rào cản kỹ thuật nghiêm trọng khi đối mặt với dữ liệu du lịch địa phương: hiện tượng ảo giác thực thể (Entity Hallucination) dẫn đến việc tự bịa đặt các địa danh không có thật, và sự thiếu hụt nhận thức không gian (Space-agnostic limitation) sinh ra các lịch trình di chuyển ziczac, vi phạm cự ly đi lại thực tế và chia cắt địa hình tự nhiên (như hai bờ sông Hương tại Cố đô Huế).

Khóa luận này đề xuất và hiện thực hóa một giải pháp tích hợp toàn diện mang tên Spatial GraphRAG, kết hợp chặt chẽ giữa Đồ thị Tri thức Không gian (Spatial Knowledge Graph gồm 9.525 nút thực thể và 90.948 cạnh quan hệ), kho dữ liệu đánh giá du khách quy mô lớn (1,69 triệu bài nhận xét từ TripAdvisor), và khung điều phối tác tử tự suy ngẫm (Agentic RAG) trên nền tảng LangGraph vận hành hoàn toàn cục bộ bằng mô hình Qwen2.5-3B trên máy trạm cá nhân (Zero API Cost). Đồng thời, đề tài xây dựng bộ từ điển ẩm thực 570 món ăn phục vụ khai phá thực đơn ngầm và phát triển giải thuật lập lịch trình phân cụm địa lý (Spatial Itinerary Planner) kết nối cổng API RESTful chuẩn mở sẵn sàng tích hợp với nền tảng công dân số Hue-S.

Kết quả kiểm thử thực nghiệm trên bộ benchmark 12 kịch bản du lịch thực tế tại Huế khẳng định mô hình Spatial GraphRAG đạt tỷ lệ bảo chứng thực thể tuyệt đối (Grounding Rate 100\%), triệt tiêu hoàn toàn tỷ lệ ảo giác và lỗi vi phạm cự ly (0\% Spatial Violation), rút ngắn khoảng cách di chuyển trung bình xuống 1,25~km (so với 4,5~km của Naive RAG và 6,8~km của Zero-shot LLM), đồng thời nâng độ chính xác món ăn đặc sản lên 96,5\% với thời gian đáp ứng thời gian thực 2,8 giây.
